\documentclass[10pt,a4paper]{amsart}
\usepackage[margin=19mm]{geometry}
\usepackage{amsmath,amssymb,mathtools}
\usepackage[scr=rsfs,cal=euler]{mathalfa}
\usepackage{xfrac}
\usepackage[unicode,hidelinks,bookmarks=false]{hyperref}
\newcommand{\ie}{i.e.\ }
\newcommand{\cf}{cf.\ }
\newcommand{\resp}{respectively\ }
\newcommand{\Subsection}{\subsection}
\providecommand{\nobreakdash}{\nobreak\hbox{-}}
\setlength{\emergencystretch}{2em}
\multlinegap0pt
\usepackage{amssymb}
\usepackage[shortlabels]{enumitem}
\usepackage{float}
\usepackage{mathtools}
\newtheorem{proposition}{Proposition}
\title{Dirac Operators, APS Boundary Conditions, and Spectral Flow on a Finite Warped Cylinder}
\author{Taro Kimura$^*$, Sanchita Sharma}
\date{Source: arXiv:2603.23275v1 (CC BY 4.0)}
\begin{document}
\maketitle

\noindent\emph{Source and licence.} Dirac Operators, APS Boundary Conditions, and Spectral Flow on a Finite Warped Cylinder. Version arXiv:2603.23275v1 (CC BY 4.0), distributed by arXiv under the Creative Commons Attribution 4.0 International licence, http://creativecommons.org/licenses/by/4.0/, as recorded in the arXiv licence metadata for this exact version and shown on the abstract page https://arxiv.org/abs/2603.23275v1. Full licence notice: \texttt{licenses/arxiv-2603.23275-CC-BY-4.0.txt}.\par\smallskip

\noindent\textit{Editorial scope.} Counted unit: the article's proposition prop:explicit-zero-mode-criterion (source0.tex lines 1661--1684). The article's complete original proof of it is delivered with it (source0.tex lines 1686--1724, one \texttt{proof} environment of 39 lines). No mathematics is written, repaired or re-derived: the statement and the proof are the article's own lines, in the article's own notation, and no retained line is substituted or re-flowed.  Retained uncounted as the source-local context the proof needs: lines 1634--1638; lines 1653--1658.  Nothing was omitted from any retained span, so the package carries no omission record.  No retained line contains a citation, so the wrapper carries no bibliography and no bibliography entry.  No retained line contains a figure environment, an image-inclusion command or drawing code, so no image is delivered and the package declares no asset dependency.  Editorial formatting only, in addition: an AMS \texttt{amsart} preamble that restates the article's own declarations and macros used by the retained lines, namely the environments \texttt{proposition} on separate counters, together with the standard packages those lines need.  The retained environments print in this document as Proposition 1 in order of appearance, whereas the article prints the same items with the numbering of its own full document.  The complete native source of the article as distributed in the arXiv e-print for this exact version is bundled unchanged as \texttt{sources/197069/source0.tex}.
\begin{equation}\label{eq:non-degenerate-delta}
\delta \neq \frac{2}{\ell(T)},
\qquad
\ell(T)=\int_0^T \frac{d\tau}{f(\tau)}.
\end{equation}

\begin{equation}\label{eq:F_matching_clean}
F(s,\lambda,k)=
(1+\alpha_T(s,k))\,u(T;s,\lambda,k)
+
(1-\alpha_T(s,k))\,w(T;s,\lambda,k).
\end{equation}

\begin{proposition}[Explicit zero-mode criterion for the regularized family]
\label{prop:explicit-zero-mode-criterion}
Fix a mode \(k\) and write \(m(s,k)=k+A(s)\). For \(\lambda=0\), the matching function \eqref{eq:F_matching_clean} is
\begin{equation}\label{eq:F-zero-explicit}
F(s,0,k)
=
\Bigl(1-\tanh\!\frac{m(s,k)}{\delta}\Bigr)^2 e^{\,m(s,k)\ell(T)}
-
\Bigl(1+\tanh\!\frac{m(s,k)}{\delta}\Bigr)^2 e^{-\,m(s,k)\ell(T)}.
\end{equation}
Equivalently,
\begin{equation}\label{eq:F-zero-factor}
F(s,0,k)=0
\quad\Longleftrightarrow\quad
m(s,k)\Bigl(\ell(T)-\frac{2}{\delta}\Bigr)=0.
\end{equation}
In particular, under \eqref{eq:non-degenerate-delta},
\begin{equation}
F(s,0,k)=0
\quad\Longleftrightarrow\quad
m(s,k)=0.
\end{equation}
Hence, zero modes of the regularized family occur exactly at the boundary zeros.
\end{proposition}

\begin{proof}
At \(\lambda=0\), the system decouples:
\begin{equation}
u(t)=u(0)e^{\,m(s,k)\ell(t)},
\qquad
w(t)=w(0)e^{-\,m(s,k)\ell(t)},
\end{equation}
where
\begin{equation}
\ell(t)=\int_0^t \frac{d\tau}{f(\tau)}.
\end{equation}
With
\begin{equation}
x(0;s,0,k)=v_0(s,k)
=
\binom{1-\alpha_0(s,k)}{-(1+\alpha_0(s,k))},
\qquad
\alpha_0(s,k)=\tanh\!\Bigl(\frac{m(s,k)}{\delta}\Bigr),
\end{equation}
we get
\begin{equation}
u(T)=(1-\alpha_0(s,k))e^{\,m(s,k)\ell(T)},
\qquad
w(T)=-(1+\alpha_0(s,k))e^{-\,m(s,k)\ell(T)}.
\end{equation}
Since \(\alpha_T(s,k)=-\alpha_0(s,k)\), substitution into \eqref{eq:F_matching_clean} gives \eqref{eq:F-zero-explicit}. Using
\begin{equation}
\frac{1+\tanh z}{1-\tanh z}=e^{2z},
\end{equation}
the equation \(F(s,0,k)=0\) is equivalent to
\begin{equation}
e^{2m(s,k)\ell(T)}
=
\left(\frac{1+\tanh(m(s,k)/\delta)}{1-\tanh(m(s,k)/\delta)}\right)^2
=
e^{4m(s,k)/\delta},
\end{equation}
which is exactly \eqref{eq:F-zero-factor}. The last statement follows from \eqref{eq:non-degenerate-delta}.
\end{proof}

\end{document}
